\chapter{Conclusiones y Trabajo Futuro}
\section{Conclusiones}
El cierre del Sprint 2 transforma el prototipo de cuentas en TransDrakkar, una aplicación de transporte de mobiliario con perfiles de cliente, administrador y conductor. El incremento conecta solicitud, revisión, planificación y entrega, incluyendo incidencias y consulta del historial. Se mantiene el desarrollo incremental: la base de acceso del Sprint 1 se conserva y el segundo sprint añade las operaciones logísticas.

La separación entre presentación, comunicación, API, lógica y persistencia permite concentrar reglas y comprobar permisos en el servidor. Las transacciones, índices de reserva, revisiones y claves idempotentes resuelven problemas de consistencia, concurrencia y reintentos. El historial permite explicar el recorrido de cada solicitud y mantener el seguimiento autorizado por perfil.

La integración mediante ramas y PR aporta trazabilidad y comprobaciones previas al despliegue. La revisión del PR 42 supera 85 casos en CI, con MongoDB real y contenedor, y se complementa con una ejecución visual sobre datos ficticios. El incremento está desplegado en Cloud Run. El resultado técnico se distingue del seguimiento: DoneTonic indica 100\% completado y cero pendiente, pero sus gráficos presentan discrepancias con las estimaciones de los PBI transcritos, que se hacen explícitas en el anexo.

El objetivo de desarrollar y documentar un producto por incrementos se concreta en dos sprints con requisitos, código, pruebas y evidencias. El alcance de transporte del Sprint 2 está implementado y validado; ello no implica certificar integraciones reales de correo/Google ni funcionalidades ajenas al alcance. Los registros aportados suman 26,8 horas para los ocho PBI del Sprint 2; no se comparan con la estimación como si su unidad estuviera confirmada.

\section{Trabajo Futuro}
\begin{itemize}
\item Ampliar la gestión operativa con edición o cancelación de rutas y correcciones administrativas, definiendo previamente sus estados, permisos y efectos sobre reservas e historial.
\item Incorporar mapas u optimización de recorridos y control de capacidad solo tras acordar unidades y medidas verificables; no convertir automáticamente bultos a peso o volumen sin datos suficientes.
\item Completar las evidencias reales de correo y Google, comprobar reinicio y caducidad en el entorno desplegado y revisar los límites por proceso antes de ampliar instancias. Mantener alineados los agregados de DoneTonic, los PBI y el registro de dedicación.
\end{itemize}
